\usepackage[automark]{scrlayer-scrpage}
\pagestyle{scrheadings}
\usepackage{xcolor}
\usepackage{geometry}
\usepackage{graphicx}
\usepackage{pdfpages}
\usepackage{tcolorbox}

% USAR FUENTES CLÁSICAS DE LaTeX (sin fontspec para evitar problemas)
\usepackage{mathptmx}  % Times-like font
\usepackage[T1]{fontenc}
\usepackage[utf8]{inputenc}

% NUMERACIÓN CONSECUTIVA DE FIGURAS Y TABLAS
\usepackage{chngcntr}
\counterwithout{figure}{chapter}
\counterwithout{table}{chapter}

% COLORES BASADOS EN LA TAPA (verde pálido)
\definecolor{chaptercolor}{RGB}{85, 180, 120}    % Verde pálido similar a la tapa
\definecolor{accentcolor}{RGB}{242, 169, 0}      % Dorado como acento
\definecolor{lightgreen}{RGB}{200, 230, 210}     % Verde muy pálido para fondos

% TAPA AL INICIO
\AtBeginDocument{%
  \includepdf[pages=-]{TAPA.pdf}
}

% CONTRATAPA AL FINAL
\AtEndDocument{%
  \includepdf[pages=-]{CONTRATAPA.pdf}
}

% Comando para DOI del capítulo (forzando espaciado pequeño)
\newcommand{\chapterdoi}[1]{%
  \vspace*{-0.5em}  % Usar \vspace* para forzar
  \begin{center}
    \textcolor{chaptercolor}{\rule{0.3\textwidth}{0.5pt}}\\[-0.2em]  % Espaciado negativo
    \small\textbf{\color{chaptercolor}DOI:} \textrm{#1}\\[-0.2em]    % Espaciado negativo
    \textcolor{chaptercolor}{\rule{0.3\textwidth}{0.5pt}}
  \end{center}
  \vspace*{-0.3em}  % Espaciado negativo y forzado
}

% Comando personalizado para autores (forzando espaciado pequeño)
\newcommand{\chapterauthors}[1]{%
  \vspace*{-0.5em}  % Reducir espacio antes
  \begin{center}
    \Large\textit{\color{chaptercolor}#1}
  \end{center}
  \vspace*{-0.5em}  % Reducir espacio después
}

% Comando para afiliaciones (forzando espaciado pequeño)
\newcommand{\chapteraffiliations}[1]{%
  \vspace*{-0.3em}  % Reducir espacio antes
  \begin{center}
    \small\color{gray}#1
  \end{center}
  \vspace*{-0.2em}  % Espaciado muy pequeño
  \noindent\textcolor{chaptercolor}{\rule{\textwidth}{0.8pt}}
  \vspace*{0.2em}   % Espaciado mínimo después de línea
  \newpage
}

% Comando mejorado para información de contribuciones CRediT (sin sombra)
\newcommand{\creditcontributions}[1]{%
  \vspace{1em}
  \begin{tcolorbox}[
    colback=white,
    colframe=chaptercolor,
    boxrule=1.5pt,
    arc=8pt,
    left=15pt,
    right=15pt,
    top=12pt,
    bottom=12pt
  ]
    \small
    \textbf{\color{chaptercolor}Contribuciones de autoría (CRediT):}\\[0.8em]
    \color{black}#1\\[0.8em]
    \footnotesize\textit{\color{gray}Taxonomía CRediT (Contributor Roles Taxonomy):}\\[0.3em]
    \footnotesize\color{gray}• Versión original (inglés): https://credit.niso.org/\\
    \footnotesize\color{gray}• Traducción al español: https://terminologias.filo.uba.ar/vocabs/1007
  \end{tcolorbox}
  \vspace{0.5em}
}

% Comando mejorado para disponibilidad de datos (sin sombra)
\newcommand{\dataavailability}[1]{%
  \vspace{1em}
  \begin{tcolorbox}[
    colback=white,
    colframe=chaptercolor,
    boxrule=1.5pt,
    arc=8pt,
    left=15pt,
    right=15pt,
    top=12pt,
    bottom=12pt
  ]
    \small
    \textbf{\color{chaptercolor}Declaración de disponibilidad de datos}\\[0.8em]
    \color{black}#1
  \end{tcolorbox}
  \vspace{1em}
  \newpage  
}

% Mejorar el estilo del título de capítulo
\addtokomafont{chapter}{\color{chaptercolor}}

% Agregar línea decorativa después del título
\renewcommand*{\chapterheadendvskip}{%
  \vspace{0.5em}
  \noindent\textcolor{chaptercolor}{\rule{\textwidth}{1pt}}
  \vspace{1em}
}

% En header.tex, agregar después de los otros comandos

% Redefinir la página de título para incluir el logo
\renewcommand*{\maketitle}{%
  \begin{titlepage}
    \centering
    
    % Logo de la editorial (ajusta el tamaño según necesites)
    \includegraphics[width=4cm]{logo-paideia.png}\\[2cm]
    
    % Título
    {\huge\bfseries Editar, publicar y financiar ciencia en América Latina}\\[0.5cm]
    
    % Subtítulo
    {\Large\textit{Perspectivas, experiencias y distopías en las dinámicas de la comunicación científica en la región}}\\[2cm]
    
    % Autores/Coordinadores
    {\large\textbf{Coordinadores:}}\\[0.5cm]
    {\large Gimena del Rio Riande}\\
    {\large Patricio Iván Pantaleo}\\[2cm]
    
    % Fecha
    {\large A COMPLETAR}\\[2cm]
    
    \vfill
    
    % Editorial al pie
    {\large\textbf{Paideia Editorial}}\\
    {\normalsize Córdoba, Argentina}
    
  \end{titlepage}
  \cleardoublepage
}